\subsection{Recommendation Process Overview}
The reviewer recommendation system uses a multi-step process to identify suitable reviewers while avoiding conflicts of interest:

\begin{enumerate}
    \item Query potential reviewers by publication count (expertise proxy)
    \item Apply conflict filtering (coauthorship, institutional)
    \item Collect metadata for scoring and diversity analysis
    \item Apply a multi-component scoring algorithm
    \item Use a hybrid selection strategy balancing score and diversity
    \item Return top-N candidates with detailed metadata
\end{enumerate}

\subsection{Conflict Filtering}
The system identifies and excludes three types of conflicts:

\begin{itemize}
    \item \textbf{Recent co-authorship:} Reviewers who co-authored papers with submission authors within a configurable time frame
    \item \textbf{Frequent co-authorship:} Reviewers with more than a threshold number of co-authored papers with submission authors
    \item \textbf{Institutional conflicts:} Reviewers affiliated with the same institutions as submission authors
\end{itemize}

\subsection{Scoring Methodology}
For each eligible reviewer, the algorithm calculates a composite score with the following components:

\begin{enumerate}
    \item \textbf{Base Expertise Score (0-5 points):} Based on publication count
    \begin{verbatim}
    // Base score from publication count (experience)
    let score = Math.min(reviewer.publication_count / 10, 5);
    \end{verbatim}
    
    \item \textbf{Venue Diversity Score:} Rewards reviewers who publish in venues that are uncommon among other potential reviewers
    \begin{verbatim}
    // Venues that are less common get higher scores
    const venueRarity = 1 / (venueFrequency[venue.journal] || 1);
    venueDiversityScore += venueRarity * 2;
    \end{verbatim}
    
    \item \textbf{Institutional Diversity Score:} Rewards reviewers from institutions that are less represented
    \begin{verbatim}
    // Less common institutions get higher scores
    const instRarity = 1 / (institutionFrequency[inst.institution_name] || 1);
    institutionDiversityScore += instRarity * 2;
    \end{verbatim}
    
    \item \textbf{Experience Score (0-3 points):} Based on years active in publishing, with square root applied for diminishing returns
    \begin{verbatim}
    const yearsActive = reviewer.years_active.last_year - 
                        reviewer.years_active.first_year + 1;
    score += Math.min(Math.sqrt(yearsActive) / 2, 3);
    \end{verbatim}
    
    \item \textbf{Collaboration Network Score (0-2 points):} Based on number of co-authors, with logarithmic scaling
    \begin{verbatim}
    // Broader networks indicate greater field awareness
    score += Math.min(Math.log10(reviewer.coauthor_count || 1), 2);
    \end{verbatim}
\end{enumerate}

\subsection{Key Metadata Queries}
The system utilizes several specialized queries to gather metadata for scoring:

\begin{verbatim}
/* Publication venues query */
SELECT DISTINCT p.journal, COUNT(*) as count
FROM publication_authors pa
JOIN publications p ON pa.publication_key = p.key
WHERE pa.person_key = $1 AND p.journal IS NOT NULL
GROUP BY p.journal
ORDER BY count DESC
LIMIT 3

/* Years of activity query */
SELECT MIN(p.year) as first_year, MAX(p.year) as last_year
FROM publication_authors pa
JOIN publications p ON pa.publication_key = p.key
WHERE pa.person_key = $1 AND p.year IS NOT NULL

/* Co-authors count query */
SELECT COUNT(DISTINCT pa2.person_key) as coauthor_count
FROM publication_authors pa1
JOIN publication_authors pa2 ON pa1.publication_key = pa2.publication_key 
                           AND pa1.person_key != pa2.person_key
WHERE pa1.person_key = $1

/* Institutional affiliations query */
SELECT institution_name, city, country, COUNT(*) as frequency
FROM author_affiliations
WHERE author_key = $1
GROUP BY institution_name, city, country
ORDER BY frequency DESC
LIMIT 3
\end{verbatim}

\subsection{Diversity-Aware Selection Strategy}
Rather than simply selecting the top N scored reviewers, the algorithm uses a hybrid strategy:

\begin{enumerate}
    \item First, select the highest-scoring reviewer
    \item Then, alternate between selecting:
    \begin{itemize}
        \item Reviewers who add diversity to the pool (on odd selections)
        \item The next highest-scoring reviewer (on even selections)
    \end{itemize}
\end{enumerate}

The diversity check evaluates whether a reviewer publishes in venues or belongs to institutions not already covered by selected reviewers:

\begin{verbatim}
const addsDiversity = (reviewer) => {
    // If we have fewer than 2 reviewers, any will do
    if (selectedReviewers.length < 2) return true;
    
    // Check for new venues
    for (const venue of reviewer.venues) {
        if (!selectedVenues.has(venue.journal)) {
            return true;
        }
    }
    
    // Check for new institutions
    if (reviewer.institutions) {
        for (const inst of reviewer.institutions) {
            if (!selectedInstitutions.has(inst.institution_name)) {
                return true;
            }
        }
    }
    
    return false;
};
\end{verbatim}

This approach ensures a balance between selecting the most qualified reviewers and maintaining diversity in research focus and institutional backgrounds. 